\documentclass[english, parskip=full]{scrartcl}
\usepackage[utf8]{inputenc}  % Kodierung der Datei
\usepackage[ngerman]{babel}  % Sprache auf Deutsch 
\usepackage{color}
\usepackage{graphicx, gensymb}
\usepackage{amsfonts, cancel, amssymb}
\usepackage{amsmath, amsthm, array, esint}
\usepackage{booktabs, multirow}  % schönere Tabellen
\usepackage{caption}  % Anpassung der Captions
\usepackage{scrlayer-scrpage}    % Manipulation des Seitenstils
\pagestyle{scrheadings}  % Stil für die Seite setzen
\automark{section}  % Abschnittsnamen als Seitenbeschriftung 
\setheadsepline{.5pt}    % Kopzeile durch Linie abtrennen
\usepackage[hidelinks]{hyperref}  % Links und weitere PDF-Features
\usepackage{typearea, tikz, pgffor, verbatim}
\usetikzlibrary{calc, arrows.meta,arrows, graphs, quotes, positioning}
\usepackage[cache=false, outputdir=build]{minted}
\usemintedstyle{vs}
\usepackage{placeins} % provides \FloatBarrier

\definecolor{bg}{rgb}{0.9,0.9,0.9}
\newcommand\numberthis{\addtocounter{equation}{1}\tag{\theequation}}
\usepackage{svg}
\usepackage{siunitx}
\usepackage{physics}
\usepackage{tensor}
\renewcommand{\bra}{}
\newcommand{\kom}{}
\newcommand{\brak}{}
\renewcommand{\braket}{}
\newcommand{\intx}{}
\newcommand{\intr}{}
\newcommand{\kla}{}
\newcommand{\klam}{}
\newcommand{\klammer}{}
\newcommand{\oper}{}
\newcommand{\tecxt}{}
\newcommand{\Bra}{}
\newcommand{\Ket}{}
\renewcommand{\i}{\text{i}}
\newcommand{\e}{\text{e}}
\newcommand*{\Fourier}{\textsc{Fourier}\ }
\newcommand*{\F}{\mathcal{F}}
\newcommand*{\sinc}{\text{sinc\,}}
\newcommand{\conv}{\mathop{\scalebox{3.5}{\raisebox{-0.38ex}{\makebox[0.5\width][c]{$\ast$}}}}\limits}

\newcommand*{\titel}{Adiabatische Näherung}
\newcommand*{\autor}{Joachim Schwardt}

\hypersetup{pdfauthor={\autor}, pdftitle={\titel}}

%\titlehead{titlehead}
\subject{Theoretische Festkörperphysik}
\title{\titel}
\author{\autor}
\date{\today}

\begin{document}

%\begin{titlepage}
\maketitle  % Titel setzen
%\tableofcontents  % Inhaltsverzeichnis setzen
%\end{titlepage}

\section{Zwei gekoppelte Oszillatoren}
\subsection*{a)}
Gegeben ist der Hamilton-Operator
\begin{align}
H &= \frac{p_1^2}{2M} + \frac{p_2^2}{2m} + \frac{K}{2}x_1^2 + \frac{k}{2} \kla\left( x_1 - x_2 \right)^2.
\end{align}
Sei $M_0 = 0$ und $M_1 = m$.
Wir transformieren unsere Impulse und Koordinaten gemäß $p_\mu \rightarrow \tilde{p}_\mu := \sqrt{M_\mu}p_\mu$ und $x_\mu \rightarrow \tilde{x}_\mu := \frac{1}{\sqrt{M_\mu} } x_\mu$.
Damit lautet der Hamilton-Operator dann
\begin{align}
H &= \frac{1}{2}\sum_{\mu} \tilde{p}_\mu^2 + \frac{1}{2}\sum_{\mu\nu} D_{\mu\nu} \tilde{x}_\mu \tilde{x}_\nu.
\end{align}
In der allgemeinen Formulierung der Vorlesung ist die dynamische Matrix als
\begin{align}
D_{n\mu j}^{n'\mu'j'} := \frac{1}{\sqrt{M_\mu M_\nu}} \frac{\partial^2V}{\partial R_{n\mu j} \partial R_{n'\mu' j'}}
\end{align}
definiert.
Hier haben wir allerdings nur eine ``Einheitszelle'' (also $n\equiv 1$) und auch nur eine Raumdimension (also $j\equiv 1$).
Damit vereinfacht sich die Formulierung zu
\begin{align}
D_{\mu\nu} = \frac{1}{\sqrt{M_\mu M_\nu}} \frac{\partial^2}{\partial \tilde{x}_\mu \partial \tilde{x}_\nu} \kla\left( \frac{K}{2} \tilde{x}_1^2 + \frac{k}{2}(\tilde{x}_1 - \tilde{x}_2)^2 \right) = \begin{pmatrix}
\frac{K+k}{M} & - \frac{k}{\sqrt{mM}} \\ - \frac{k}{\sqrt{mM}} & \frac{k}{m}
\end{pmatrix}.
\end{align}
Um diese Matrix zu diagonalisieren, bestimmen wir zunächst die Eigenwerte $\lambda$ aus der Gleichung
\begin{align}
0 &= \lambda^2 - \tecxt\text{tr\,}D + \det D = \lambda^2 - \kla\left( \frac{K+k}{M} + \frac{k}{m} \right) \lambda + \frac{kK}{mM}.
\end{align}
Die Lösungen lauten 
\begin{align}
\lambda &= \frac{K+k}{2M} + \frac{k}{2m} \pm \sqrt{\frac{k^2}{4m^2} + \frac{K^2 + 2kK + k^2}{4M^2} + \frac{kK+k^2}{2mM} - \frac{kK}{mM}} = \\
&= \frac{K+k}{2M} + \frac{k}{2m} \pm \sqrt{\frac{K^2}{4M^2} + \frac{kK}{2mM} \kla\left( \frac{m}{2M} - 1 \right) + \frac{k^2}{4m^2} \kla\left( 1 + \frac{m}{M} \right)^2}.
\end{align}
Wir identifizieren die Eigenwerte als $\omega_{1,2}^2$.
Das System kann nun als zwei voneinander unabhängige harmonische Oszillatoren mit den entsprechenden Frequenzen interpretiert werden.
Die Gesamtenergie ergibt sich daher zu
\begin{align}
E_{\nu_1\nu_2} &= \omega_1 \kla\left( \nu_1 + \frac{1}{2} \right) + \omega_2 \kla\left( \nu_2 + \frac{1}{2} \right).
\end{align}

%Die Eigenwertgleichung ist $D\vec{u} = \lambda \vec{u}$, oder explizit
%\begin{align}
%\begin{pmatrix}
%(K+k)u_1 - ku_2 \\ -ku_1 + ku_2
%\end{pmatrix} &= \lambda \begin{pmatrix}
%u_1 \\ u_2
%\end{pmatrix}.
%\end{align}
%Als Lösungen ergeben sich $u_2 = \frac{1}{k}\klam\left[ K+k-\lambda_+ \right]u_1$ und $u_1 = \frac{1}{k} \klam\left[ k - \lambda_- \right]u_2$, also die Eigenvektoren
%\begin{align}
%\vec{u}_1 &= \frac{1}{N} \begin{pmatrix}
%k \\ K+k-\lambda_+
%\end{pmatrix} = \frac{1}{N} \begin{pmatrix}
%k \\ \frac{K}{2} - \kappa
%\end{pmatrix} = \frac{1}{\sqrt{2\kappa (\kappa- K)}} \begin{pmatrix}
%k \\ \frac{K}{2} - \kappa
%\end{pmatrix},\\
%\vec{u}_2 &= \frac{1}{N} \begin{pmatrix}
%k - \lambda_- \\ k
%\end{pmatrix} = \frac{1}{N} \begin{pmatrix}
%\kappa - \frac{K}{2} \\ k
%\end{pmatrix} = \frac{1}{\sqrt{2\kappa (\kappa - K)}} \begin{pmatrix}
%\kappa - \frac{K}{2} \\ k
%\end{pmatrix}.
%\end{align}
%Mit der Transformationsmatrix $U:= (\vec{u}_1, \vec{u}_2)$ können wir nun neue Koordinaten gemäß
%\begin{align}
%\vec{y} &:= \begin{pmatrix}
%y_1 \\ y_2
%\end{pmatrix} = U\begin{pmatrix}
%x_1 \\ x_2
%\end{pmatrix} = \frac{1}{N}\begin{pmatrix}
%kx_1 + \kla\left( \frac{K}{2} - \kappa \right)x_2 \\ 
%\kla\left( \kappa - \frac{K}{2} \right)x_1 + kx_2
%\end{pmatrix}
%\end{align}
%einführen.
%Insbesondere gilt
%\begin{align}
%\frac{\partial}{\partial x_1} &= \frac{\partial y_1}{\partial x_1} \frac{\partial}{\partial y_1} +  \frac{\partial y_2}{\partial x_1} \frac{\partial}{\partial y_2} = k \frac{\partial}{\partial y_1} +  \kla\left( \kappa - \frac{K}{2} \right) \frac{\partial}{\partial y_2}.
%\end{align}

\subsection*{b)}
In der adiabatischen Näherung fassen wir $T_a := \frac{p_1^2}{2M}$ als die ``atomare'' kinetische Energie auf und vernachlässigen diese zunächst.
Übrig bleibt das ``elektronische'' Problem
\begin{align}
H_0 &= \frac{p_2^2}{2m} + \frac{K}{2}x_1^2 + \frac{k}{2}(x_1 - x_2)^2.
\end{align}
Dieses kann mit $y_2 := x_2 - x_1$ sehr einfach gelöst werden.
Es ergibt sich
\begin{align}
H_0 &= \frac{p_2^2}{2m} + \frac{K}{2}x_1^2 + \frac{k}{2}y_2^2,
\end{align}
und damit die Energien $\epsilon_{\nu_2}(x_1) = \sqrt{\frac{k}{m}}\kla\left( \nu_2 + \frac{1}{2} \right) + \frac{K}{2}x_1^2$.
\\
Jetzt lösen wir noch das verbliebene ``atomare'' Problem
\begin{align}
H_0 &= \frac{p_1^2}{2M} + \epsilon_{\nu_1}(x_1) = \frac{p_1^2}{2M} + \frac{K}{2}x_1^2 + \sqrt{\frac{k}{m}}\kla\left( \nu + \frac{1}{2} \right).
\end{align}
Die Lösung führt auf die Gesamtenergien
\begin{align}
E_{\nu_1\nu_2} &= \sqrt{\frac{K}{M}}\kla\left( \nu_1 + \frac{1}{2} \right) + \sqrt{\frac{k}{m}}\kla\left( \nu_2 + \frac{1}{2} \right).
\end{align}
Das entspricht gerade der Energie zweier nicht-gekoppelter harmonischer Oszillatoren.

\subsection*{c)}
Die beiden Lösungen haben bereits dieselbe Struktur, allerdings unterscheiden sich die Frequenzen der zugrunde liegenden Oszillatoren.
Wir betrachten die exakten Ergebnisse für $\frac{m}{M} \ll 1$
\begin{align}
\omega_{1,2}^2 &= \frac{1}{m} \kla\left( \frac{k}{2} + [K+k]\frac{m}{2M} \pm \sqrt{\frac{K^2m^2}{4M^2} + \frac{kKm}{2M} \kla\left( \frac{m}{2M} - 1 \right) + \frac{k^2}{4} \kla\left( 1 + \frac{m}{M} \right)^2} \right) \simeq \\
&\simeq \frac{1}{m} \kla\left( \frac{k}{2} + [K+k]\frac{m}{2M} \pm \sqrt{\frac{k^2}{4} + \frac{k^2m}{2M} - \frac{kKm}{2M}} \right) \simeq \\
&\simeq \frac{1}{m} \kla\left( \frac{k}{2} + [K+k]\frac{m}{2M} \pm \frac{k}{2} \klam\left[ 1 + \frac{m}{M} - \frac{Km}{kM}\right] \right).
\end{align}
Für die unterschiedlichen Vorzeichen ergeben sich damit
\begin{align}
\omega_1^2 &= \frac{1}{m} \kla\left( k + [K+k]\frac{m}{2M} + \frac{km}{2M} - \frac{Km}{2M} \right) = \frac{k}{m} \kla\left( 1  + \frac{m}{M} \right),\\
\omega_2^2 &= \frac{1}{m} \kla\left( [K+k]\frac{m}{2M} - \frac{km}{2M} + \frac{Km}{2M} \right) = \frac{K}{M}.
\end{align}
Das sind (für $\frac{m}{M} \rightarrow 0$) gerade die Frequenzen aus der adiabatischen Näherung.

%\newpage
\section{Graphen}

\begin{figure}[!ht]
\centering
\begin{tikzpicture}   
    \pgfmathsetmacro{\a}{1.0}
    \pgfmathsetmacro{\d}{sqrt(3.0)*\a}
    \pgfmathsetmacro{\h}{1.5*\a}
    \pgfmathsetmacro{\phi}{60}
    \pgfmathsetmacro{\r}{0.05}
    \pgfmathsetmacro{\xmin}{-5}
    \pgfmathsetmacro{\xmax}{5}
    \pgfmathsetmacro{\ymin}{-2.5}
    \pgfmathsetmacro{\ymax}{4}
    
    
    \coordinate (o) at (0,0);
    \coordinate (xmax) at (\xmax,0);
    \coordinate (xmin) at (\xmin,0);
    \coordinate (ymax) at (0,\ymax);
    \coordinate (ymin) at (0,\ymin);
    
    \node[anchor=north] at (xmax) {$x$};
    \node[anchor=west] at (ymax) {$y$};
    
%    color=black!35!white
%    \draw[thick, -{Stealth[length=1.5mm,width=1.5mm]}] (kmin) arc (180:45:\R);
    
    \draw[color=black, -latex] (xmin) to (xmax);
    \draw[color=black, -latex] (ymin) to (ymax);
    
    % primitive Einheitszelle
    \filldraw[fill=green!40!white, draw=black] (o) -- (0.5*\d,\h) -- (1.5*\d,\h) -- (\d,0) -- cycle;
    
    % Wigner-Seitz-Zelle
    \draw (-0.5*\d,-0.5*\a) --++ (0,2*\h) --++ (-1.5*\d,-1.5*\a) -- cycle;
    \draw (-1.5*\d,-0.5*\a) --++ (0,2*\h) --++ (1.5*\d,-1.5*\a) -- cycle;
    \filldraw[fill=black!15!white, draw=black] (-0.5*\d,0.5*\a) --++ (0,\a) --++ (-0.5*\d,0.5*\a) --++ (-0.5*\d,-0.5*\a) --++ (0,-\a) --++ (0.5*\d,-0.5*\a) -- cycle;
    
    \draw[very thick, color=red!60!black, -{Stealth[length=1.5mm,width=1.5mm]}] (o) to (0,-\a);
    \node[left=0.5mm] at (0,-\a) {$\color{red!60!black}{-\vec{a}_1'}$};
    
    % Grid zeichnen
    \foreach \i in {-2,...,2}{
        \foreach \j in {0,2}{
            \draw[fill] (\i*\d,\j*\h) circle (\r);
            \draw[fill] (\i*\d + 0.5*\d,\j*\h - \h) circle (\r);
        }
    }
    \foreach \i in {-2,...,2}{
        \foreach \j in {0,2}{
            \draw[fill] (\i*\d + 0.5*\d,\j*\h - 0.5*\a) circle (\r);
            \draw[fill] (\i*\d,\j*\h - \h - 0.5*\a) circle (\r);
%            \draw[color=red] (\i*\d + 0.5*\d,\j*\h - 0.5*\a) circle (\a);
%            \draw[color=red] (\i*\d,\j*\h - \h - 0.5*\a) circle (\a);
        }
    }
    
    % Kohlenstoff-Sorten-Labels A und B
    \node[below left=-0.08cm and -0.04cm] at (o) {A};
    \node[below left=-0.08cm and -0.04cm] at (-\d,0) {A};
    \node[left] at (-0.5*\d,\h) {A};
    \node[left] at (-1.5*\d,\h) {A};
    
    \node[left] at (0,\a) {B};
    \node[left] at (-\d,\a) {B};
    \node[left] at (-0.5*\d,-0.5*\a) {B};
    \node[left] at (-1.5*\d,-0.5*\a) {B};
    
    % Wrong basis vectors
    \draw[very thick, color=blue, -{Stealth[length=1.5mm,width=1.5mm]}] (o) to (0,\a);
    \node[above left=0.6mm and 0.5mm] at (o) {$\color{blue}{\vec{a}_1'}$};
    
    \draw[very thick, color=blue, -{Stealth[length=1.5mm,width=1.5mm]}] (o) to (0.5*\d,-0.5*\a);
    \node[below right=0.3cm and 0.9mm] at (o) {$\color{blue}{\vec{a}_2'}$};
    
    % True basis vecotrs
    \draw[very thick, color=green!60!black, -{Stealth[length=1.5mm,width=1.5mm]}] (o) to (\d,0);
    \node[below right=0.0mm and 1.2cm] at (o) {$\color{green!60!black}{\vec{a}_1}$};
    
    \draw[very thick, color=green!60!black, -{Stealth[length=1.5mm,width=1.5mm]}] (o) to (0.5*\d,\h);
    \node[above right=0.4cm and 0.2cm] at (0,\a) {$\color{green!60!black}{\vec{a}_2}$};
    
%    \draw (0,0) circle (\a);
%    \draw (0,0) arc (60:0:\d);
    
\end{tikzpicture}
\caption{Skizze der Gitterstruktur von Graphen mit Strukturgröße $a$, also $|\vec{a}_{1,2}'| = a$ und $|\vec{a}_{1,2}| = \sqrt{3}a$.
Das zugrundeliegende Bravais-Gitter ist dreieckig und ist hier an den Kohlenstoff-Atomen der Sorte A orientiert.
Die primitiven Gittervektoren $\vec{a}_{1,2}$ spannen die primitive Einheitszelle, welche hier grün hinterlegt ist.
Die Wigner-Seitz-Zelle ist grau hinterlegt, wobei die Linien zur Konstruktion explizit dargestellt sind.}
\label{figure:tikz graphen}
\end{figure}

\subsection*{a)}
Beim zweidimensionalen hexagonalen Punktgitter gibt es keine ``Löcher''. 
Genauer, es müsste bei Graphen auch Gitterpunkte bei $-\vec{a}_1'$ geben (rot eingezeichnet in \autoref{figure:tikz graphen}).
Deshalb ist das Honigwabengitter allerdings auch kein Bravais-Gitter, da nicht jede Linearkombination der $\vec{a}_{1,2}'$ wieder auf dem Gitter landet.

\subsection*{b)}
Die primitiven Gittervektoren sind $\vec{a}_{1,2}$ (siehe \autoref{figure:tikz graphen}).
Es liegt ein dreieckiges Bravais-Gitter mit zweiatomiger Basis vor. 
Eine mögliche primitive Einheitszelle ist in \autoref{figure:tikz graphen} grün hinterlegt.
Die Wigner-Seitz-Zelle entspricht hier einem Hexagon mit Kantenlänge $a$.
Die Fläche ist daher gleich dem sechsfachen der gleichseitigen Dreiecke, also ist 
\begin{align}
A_\tecxt\text{WS} &= 6A_\tecxt\text{Dreieck} = 6\cdot \frac{1}{2} \cdot a \cdot \frac{\sqrt{3}a}{2} = \frac{3\sqrt{3}}{2}a^2.
\end{align}

\subsection*{c)}
Um die Basisvektoren des reziproken Gitters zu bestimmen, führen wir zunächst noch $\vec{a}_3 = (0,0,a)^\top$ als dritten Gittervektor ein.
Die anderen beiden lauten dann $\vec{a}_1 = (\sqrt{3}a,0,0)^\top$ und $\vec{a}_2 = (\frac{\sqrt{3}}{2}a,\frac{3}{2}a,0)^\top$.
Dann können wir die allgemeine Formel verwenden, wobei zunächst einmal
\begin{align}
V &= \vec{a}_3\cdot \kla\left( \vec{a}_1\times\vec{a}_2 \right) = \vec{a}_3 \cdot \frac{3\sqrt{3}}{2} a^2\vec{e}_z = \frac{3\sqrt{3}}{2}a^3.
\end{align}
Damit können wir dann die Basisvektoren berechnen
\begin{align}
\vec{b}_1 &= \frac{2\pi}{V} (\vec{a}_2\times\vec{a}_3) = \frac{4\pi}{3a}  \begin{pmatrix}
\sqrt{3}/2 \\ -1/2 \\ 0
\end{pmatrix},\\
\vec{b}_2 &= \frac{2\pi}{V} (\vec{a}_3\times\vec{a}_1) = \frac{4\pi}{3a}  \begin{pmatrix}
0 \\ 1 \\ 0
\end{pmatrix},\\
\vec{b}_3 &= \frac{2\pi}{V} (\vec{a}_1\times\vec{a}_2) = \frac{2\pi}{a}  \begin{pmatrix}
0 \\ 0 \\ 1
\end{pmatrix}.
\end{align}
\begin{figure}[!ht]
\centering
\begin{tikzpicture}[scale=1]
    \pgfmathsetmacro{\a}{1.0}
    \pgfmathsetmacro{\d}{4*pi/3/\a}
    \pgfmathsetmacro{\h}{sqrt(3)/2*\d}
    \pgfmathsetmacro{\phi}{60}
    \pgfmathsetmacro{\r}{0.05}
    \pgfmathsetmacro{\xmin}{-5}
    \pgfmathsetmacro{\xmax}{5}
    \pgfmathsetmacro{\ymin}{-5}
    \pgfmathsetmacro{\ymax}{5}
    
    
    \coordinate (o) at (0,0);
    \coordinate (xmax) at (\xmax,0);
    \coordinate (xmin) at (\xmin,0);
    \coordinate (ymax) at (0,\ymax);
    \coordinate (ymin) at (0,\ymin);
    
    \node[anchor=north] at (xmax) {$x$};
    \node[anchor=west] at (ymax) {$y$};
    
%    color=black!35!white
%    \draw[thick, -{Stealth[length=1.5mm,width=1.5mm]}] (kmin) arc (180:45:\R);
    
    % primitive Einheitszelle
%    \filldraw[fill=green!40!white, draw=black] (o) -- (0.5*\d,\h) -- (1.5*\d,\h) -- (\d,0) -- cycle;
    
    % Wigner-Seitz-Zelle
    \draw (-\h,-0.5*\d) --++ (2*\h,0) --++ (-\h,1.5*\d) -- cycle;
    \draw (-\h,0.5*\d) --++ (2*\h,0) --++ (-\h,-1.5*\d) -- cycle;
    \filldraw[fill=black!15!white, draw=black] (0.5*\h,-0.25*\d) --++ (0,0.5*\d) --++ (-0.5*\h,0.25*\d) --++ (-0.5*\h,-0.25*\d) --++ (0,-0.5*\d) --++ (0.5*\h,-0.25*\d) -- cycle;
%    \draw (o) circle (0.5*\d);    
    
    \draw[color=black, -latex] (xmin) to (xmax);
    \draw[color=black, -latex] (ymin) to (ymax);
    
    % Grid zeichnen
    \foreach \i in {-1,...,1}{
        \foreach \j in {0}{
            \draw[fill] (2*\j*\h,\i*\d) circle (\r);
        }
    }
    \foreach \i in {0,1}{
        \foreach \j in {0,1}{
            \draw[fill] (-2*\j*\h + \h,\i*\d - 0.5*\d) circle (\r);
        }
    }
    
    % Reciprocal basis vecotrs
    \draw[very thick, color=green!60!black, -{Stealth[length=1.5mm,width=1.5mm]}] (o) to (\h,-0.5*\d);
    \node[below right=0.0mm and 1.2cm] at (o) {$\color{green!60!black}{\vec{b}_1}$};
    
    \draw[very thick, color=green!60!black, -{Stealth[length=1.5mm,width=1.5mm]}] (o) to (0,\d);
    \node[above right=0.4cm and 0.2cm] at (0,0.5*\d) {$\color{green!60!black}{\vec{b}_2}$};
    
%    \draw (0,0) circle (\d);
%    \draw (2*\h,0) circle (\d);
    
\end{tikzpicture}
\caption{Skizze des reziproken Gitters mit den Basisvektoren $\vec{b}_{1,2}$.
Die erste Brillouin-Zone ist grau hinterlegt und die Konstruktionslinien sind explizit dargestellt.
Die Skizze ist maßstabsgetreu relativ zu \autoref{figure:tikz graphen}.
Da $a=1$ gewählt wurde, erscheint das reziproke Gitter deutlich dünner besetzt.}
\label{figure:tikz graphen}
\end{figure}
\\
\section{Orthogonalität}
Die Gittervektoren sind gegeben durch $\textbf{R}_n = n_1\textbf{a}_1 + n_2\textbf{a}_2 + n_3\textbf{a}_3$ mit $n_{1,2,3}\in\mathbb{Z}$.
Seien $\mathbf{q},\mathbf{q}'\in$ 1.\,BZ, also $\mathbf{q} = m_1\textbf{b}_1 + m_2\textbf{b}_2 + m_3\textbf{b}_3$ mit $m_{1,2,3}\in\mathbb{Z}$.
Dann gilt
\begin{align}
\frac{1}{N}\sum_n \e^{\i (\textbf{q}' - \textbf{q})\cdot \textbf{R}_n} &= \frac{1}{N}\sum_{n_1n_2n_3} \e^{\i \sum_{j,k=1}^3 (m_j'-m_j)n_k \textbf{b}_j\cdot \textbf{a}_k} = \\
&= \frac{1}{N}\sum_{n_1n_2n_3} \prod_{j=1}^3 \e^{2\pi\i (m_j'-m_j)n_j} = \\
&= \prod_{j=1}^3 \frac{1}{N_j} \sum_{n_j=0}^{N_j-1} \e^{2\pi\i (m_j'-m_j)n_j} = \\
&= \prod_{j=1}^3 \frac{1}{N_j} \frac{1 - \e^{2\pi\i (m_j'-m_j)N_j}}{1 - \e^{2\pi\i (m_j' - m_j)}} = \prod_{j=1}^3 \delta_{m_jm_j'} \equiv \delta_{\textbf{q}\textbf{q}'}.
\end{align}

\end{document}